\usepackage{ctex}
\usepackage[table]{xcolor}   % \rowcolor 等表格着色（机械原理章节依赖）
\usepackage{geometry,graphicx,color,subcaption}
\geometry{a4paper, top=2.5cm, bottom=2.5cm, left=2.8cm, right=2.8cm}
\usepackage{amssymb,amsmath,amsthm}                             % 数学字体
\usepackage{cancel}                                             % 数学划线（\cancel）
\usepackage{booktabs}   % 提供 \toprule, \midrule, \bottomrule    % 采用 Palatino 风格字体
\usepackage{multirow}   % 表格合并行（\multirow）
\usepackage{newclude,ulem}
\definecolor{winered}{rgb}{0.5,0,0}
\definecolor{structurecolor}{RGB}{122,122,142}
\definecolor{main}{rgb}{0.5,0,0}
\definecolor{second}{RGB}{115,45,2}
\definecolor{third}{RGB}{0,80,80}
\usepackage[colorlinks,linkcolor = winered]{hyperref}           % 定义引用的颜色
\usepackage{tcolorbox}
\tcbuselibrary{skins,breakable}
\usepackage{listings}
\usepackage{xcolor}
\lstdefinestyle{pythonstyle}{
    language=Python,
    basicstyle=\ttfamily\small,
    keywordstyle=\color{blue}\bfseries,
    commentstyle=\color{green!50!black}\itshape,
    stringstyle=\color{red},
    numberstyle=\tiny\color{gray},
    numbers=left,
    numbersep=8pt,
    frame=single,
    rulecolor=\color{black!30},
    breaklines=true,
    showstringspaces=false,
    tabsize=4,
    captionpos=b,
    xleftmargin=2em,
    framexleftmargin=1.5em
}
\lstset{style=pythonstyle}


% ------------------------------------------------------------%
% 定义定理环境
\usepackage{amsthm}
\newtheoremstyle{defstyle}{3pt}{3pt}{\kaishu}{-3pt}{
	\bfseries\color{main}}{}{0.5em}{\indent 【\thmname{#1} \thmnumber{#2}】 \thmnote{(#3)}}
\newtheoremstyle{thmstyle}{3pt}{3pt}{\kaishu}{-3pt}{
	\bfseries\color{second}}{}{0.5em}{\indent【\thmname{#1} \thmnumber{#2}】 \thmnote{(#3)}}
\newtheoremstyle{prostyle}{3pt}{3pt}{\kaishu}{-3pt}{
	\bfseries\color{third}}{}{0.5em}{\indent【\thmname{#1} \thmnumber{#2}】 \thmnote{(#3)}}

\theoremstyle{thmstyle} %theorem style
\newtheorem{theorem}{定理}[chapter]
\theoremstyle{defstyle} % definition style
% 机械原理使用自定义的无编号 definition 环境（见文件末尾 NewDocumentEnvironment 定义），
% 故此处不再用 amsthm 定义带编号的 definition，避免命令冲突。
% \newtheorem{definition}[theorem]{定义}
\newtheorem{lemma}[theorem]{引理}
\newtheorem{corollary}[theorem]{推论}
\theoremstyle{prostyle} % proposition style
\newtheorem{proposition}[theorem]{命题}
\newtheorem{example}[theorem]{例题}

\renewenvironment{proof}[1][Derivations]{\par{\heiti \textbf{#1：}} \;\songti}{\qed\par}
\newenvironment{solution}{\par\underline{\textbf{解.}} \;\kaishu}{\par}
\newenvironment{remark}{\par\underline{\textbf{注.}} \;\fangsong}{\par}
\newcommand{\intro}[1]{\rightline{\parbox[t]{5cm}{\footnotesize \fangsong\quad\quad #1 }}}
% ------------------------------------------------------------%

\usepackage{enumerate}
\usepackage{enumitem}
\setlist[enumerate,1]{label=\color{structurecolor}\arabic*.}
\setlist[enumerate,2]{label=\color{structurecolor}(\arabic*).}
\setlist[enumerate,3]{label=\color{structurecolor}\Roman*.}
\setlist[enumerate,4]{label=\color{structurecolor}\Alph*.}
\usepackage{float}
\usepackage{tikz}
\usetikzlibrary{angles, quotes, arrows.meta, calc}
% 设置章形式
% ---------------------------------- %
% \setlength{\parindent}{0pt}
\usepackage{titlesec, titletoc}
\linespread{1.5}

\usepackage{fancyhdr}
\fancypagestyle{plain}{%
	\fancyhf{} % clear all header and footer fields
	\renewcommand{\headrulewidth}{0pt}
	\renewcommand{\footrulewidth}{0pt}
}

\titlecontents{chapter}[0em]{}{\large \fangsong{第 \thecontentslabel 章\quad}}{}{\hfill\contentspage}
\titlecontents{section}[2em]{}{\thecontentslabel\quad\textcolor{blue}}{}{\titlerule*{ .} \contentspage}

\titleformat{\chapter}[display]{\LARGE}
{\color{structurecolor}\centering\small \color{structurecolor}第 \zhnumber{\arabic{chapter}} \ 章 }{0.5ex}
{\color{structurecolor}{\titlerule[1pt]}\huge \kaishu \centering \bfseries}

\titleformat{\section}[frame]
{\normalfont\color{structurecolor}}
{\footnotesize \enspace \large \textcolor{structurecolor}{\S \,\thesection}\enspace}{6pt}
{\Large\filcenter \bf \kaishu }

\titlespacing*{\section}{1pc}{*7}{*2.3}[1pc]
\titleformat{\subsection}[hang]{\bfseries}{
	\large\bfseries\color{structurecolor}\thesubsection\enspace}{1pt}{%
	\color{structurecolor}\large\bfseries\filright}
\titleformat{\subsubsection}[hang]{\bfseries}{
	\normalsize\bfseries\color{structurecolor}\thesubsubsection\enspace}{1pt}{%
	\color{structurecolor}\normalsize\bfseries\filright}

% 导言区定义格式
\newtcolorbox{myTheorem}[1]{
    colframe=blue!50!black,
    colback=blue!5!white,
    title=#1,                      % 标题由用户输入
    fonttitle=\bfseries,
    boxrule=0.5mm,
    before skip=10pt,
    after skip=10pt,
    % 保持内部 \tcbline 的分隔样式与边框一致
    segmentation style={blue!50!black, line width=0.5mm},
}
\newtcolorbox{myhypothesis}[1]{
    colframe=green!50!black,      % 边框：绿色50% + 黑色50%
    colback=green!5!white,        % 背景：绿色5% + 白色95%
    title=#1,
    fonttitle=\bfseries,
    boxrule=0.5mm,
    before skip=10pt,
    after skip=10pt,
    segmentation style={green!50!black, line width=0.5mm},  % 分割线与边框同色
}
\newtcolorbox{mydefinition}[1]{
    colframe=red!50!black,      % 边框：红色50% + 黑色50%
    colback=red!5!white,        % 背景：红色5% + 白色95%
    title=#1,
    fonttitle=\bfseries,
    boxrule=0.5mm,
    before skip=10pt,
    after skip=10pt,
    segmentation style={red!50!black, line width=0.5mm},  % 分割线与边框同色
}

% 导言区添加（需要 \usepackage{tcolorbox}）
% 导言区改为：
\tcbuselibrary{skins, breakable, theorems}
\newtcolorbox{myformula}[1][]{
    colframe=brown!50!black,
    colback=brown!5!white,
    boxrule=0.5mm,
    before skip=10pt,
    after skip=10pt,
    halign=center,
    valign=center,
    sharp corners,
    #1
}
\definecolor{DarkBlue}{RGB}{0,0,139}
\definecolor{gold}{RGB}{212,175,55}

% ============================================================
%  机械原理特有设置（保留原主文件中的自定义样式）
% ============================================================
\usepackage{siunitx}
\usepackage{yhmath}
\usepackage{wallpaper}
\usepackage{xparse}

\tcbset{
  mybox/.style = {
    enhanced, breakable,
    colback = white,
    colframe = #1!50!black,
    coltitle = white,
    fonttitle = \bfseries\kaishu,
    attach boxed title to top left = {yshift = -2mm, xshift = 3mm},
    boxed title style = {
      empty,
      boxrule = 0pt,
      colback = #1,
      rounded corners = north,
    },
    left = 3mm, right = 3mm, top = 3mm, bottom = 3mm,
    before skip = 12pt, after skip = 12pt,
    drop shadow,
    sharp corners,
  }
}

% 自定义无编号 definition 环境（可选参数为标题，蓝色盒子，支持跨页）
\NewDocumentEnvironment{definition}{o}
  {%
    \IfNoValueTF{#1}
      {\begin{tcolorbox}[
          title=定义,
          colframe=blue!50!black,
          colback=blue!5!white,
          sharp corners,
          breakable,
          fonttitle=\bfseries\kaishu,
          before skip=12pt,
          after skip=12pt
        ]}
      {\begin{tcolorbox}[
          title=定义：#1,
          colframe=blue!50!black,
          colback=blue!5!white,
          sharp corners,
          breakable,
          fonttitle=\bfseries\kaishu,
          before skip=12pt,
          after skip=12pt
        ]}
  }
  {%
    \end{tcolorbox}
  }

% 为 amsthm 定理环境添加盒子外观
\tcolorboxenvironment{theorem}{mybox={second}}
\tcolorboxenvironment{lemma}{mybox={second}}
\tcolorboxenvironment{corollary}{mybox={second}}
\tcolorboxenvironment{proposition}{mybox={third}}
\tcolorboxenvironment{example}{mybox={structurecolor}}

\tikzset{
  joint/.style={circle, fill=black, minimum size=4pt, inner sep=0pt},
  link/.style={thick},
  fixed/.style={thick, double},
  arrow/.style={-{Stealth[scale=0.8]}}
}
